\documentclass[12pt]{article}
\usepackage[utf8]{inputenc}
\usepackage[T1]{fontenc}
\usepackage{hyperref}
\title{Khengkha Language Resources for K-12 Curriculum Development}
\author{Compiled by Claude Sonnet 5.5 (claude-sonnet-5-5) for Andr\'as Kornai}
\date{}
\begin{document}
\maketitle

\section{Introduction}
Khengkha (also Kheng; ISO 639-3 xkf; Glottolog khen1241) is an East Bodish Sino-Tibetan language of Bhutan. Wikipedia gives about 40,000 speakers in Zhemgang, Trongsa and Mongar districts. The Endangered Languages Project lists it as vulnerable. Joshua Project reports about 92\% lexical similarity with Bumthang. I found no school materials in it.

\section{Primary Educational Resources}
I found no primers, textbooks or school programmes for Khengkha.

\section{Linguistic and Reference Materials}
\begin{itemize}
\item Wikipedia, ``Kheng language'': an overview. \url{https://en.wikipedia.org/wiki/Kheng_language}
\item Endangered Languages Project entry for xkf: it cites George van Driem's 2007 work on the languages of South Asia. I did not verify the bibliography. \url{https://www.endangeredlanguages.com/lang/xkf}
\item Kyoto University repository PDF on Sino-Tibetan languages, with a basic vocabulary comparison that includes Khengkha. \url{https://repository.kulib.kyoto-u.ac.jp/dspace/bitstream/2433/263979/1/Sino.Tibetan.lang_4_121.pdf}
\end{itemize}

\section{Religious and Community Resources}
\begin{itemize}
\item Joshua Project profile of the Kheng: community information from a missions database. \url{https://joshuaproject.net/people_groups/12658}
\item Asia Harvest profile of the Kheng (PDF), also a missions source. \url{https://www.people-groups.asiaharvest.org/Bhutan/Kheng.pdf}
\end{itemize}

\section{Digital and Audio Resources}
I found no audio, video or corpus resources for Khengkha.

\section{Curriculum Development Context}
The Zhemgang district document ``Understanding and Appreciating Zhemgang'' notes dialect differences in vocabulary and tense formation, but it is not a grammar. Khengkha is vigorous at home, so a curriculum could draw on community speakers. It would need new materials, an agreed orthography, and coordination with Bhutan's national school system and its language policy.

\begin{itemize}
\item Zhemgang district government: ``Understanding and Appreciating Zhemgang''. \url{https://zhemgang.gov.bt/wp-content/uploads/2025/06/Understanding-and-Appreciating-Zhemgang.pdf}
\end{itemize}

\end{document}